\chapter{The project file}
\label{ch:project}

A project file names the design, the scripts and the output folder, so that you can
run \cmd{vet apply} with no arguments. The file is called \file{hdl-vet.toml}. It is
a TOML file. \cmd{vet init} writes one for you (chapter~\ref{ch:first}).

\section{The keys}

\begin{description}
\item[\code{[design]} \code{filelist}] The top-level file list. It is required when
  the project gives the design.
\item[\code{[design]} \code{top}] The top module. It is optional, but name it
  (chapter~\ref{ch:signoff}).
\item[\code{[design]} \code{base}] The base of the file list. It works as \opt{--base}.
  You can give it only together with \code{filelist}.
\item[\code{[apply]} \code{scripts}] A list of script files, in the order in which they
  run. It is optional.
\item[\code{[apply]} \code{output}] The output folder. It is optional.
\end{description}

File names in the project file start from the folder of the project file, and use
\code{/}. hdl-vet takes an absolute name as it is. Here is the file that \cmd{vet init}
wrote in chapter~\ref{ch:first}:

\begin{itemize}
\item \code{filelist = "../design/design.f"}
\item \code{top = "soc_top"}
\item \code{scripts = ["scripts/01_example.py"]}
\item \code{output = "out"}
\end{itemize}

hdl-vet reads the file strictly. These things are an error:

\begin{itemize}
\item an unknown table or key
\item a value of the wrong type
\item an empty string
\item a nested table
\item text that is not TOML
\end{itemize}

The error names the file and the key. hdl-vet guesses nothing. The file must be UTF-8.

\section{Where it is looked for}

\cmd{vet apply} uses \file{hdl-vet.toml} in the current folder, or the file that
\opt{--project} names. It does not look in parent folders. \opt{--no-project} ignores
the file. These things are an error:

\begin{itemize}
\item \opt{--project} and \opt{--no-project} together
\item an \opt{--project} file that does not exist
\item an \opt{--project} name that is a folder
\end{itemize}

Whenever hdl-vet uses a project file, the first line on the error output is
\code{project: FILE}, with the full name. A project never runs without notice.

\section{What the command line replaces}

\begin{itemize}
\item \opt{--top} replaces \code{top} of the project.
\item \opt{-f} replaces \code{filelist} and \code{base} of the project together,
  because a base belongs to its file list. \opt{--base} replaces only the base.
\item \opt{-s} adds scripts. The scripts of the project run first, then the scripts of
  the command line. When you give \opt{-f} and the project has scripts, hdl-vet prints
  this note: \code{note: the project's scripts run first: ...}.
\item hdl-vet uses \code{output} of the project only when you give none of \opt{-o},
  \opt{--save} and \opt{--dry-run}. \cmd{vet apply --save FILE} saves and writes no
  tree, as without a project.
\item With \opt{--resume}, hdl-vet does not use the design of the project. The saved
  session has its own design. hdl-vet compares the scripts of the project with the
  scripts that the session applied (chapter~\ref{ch:session}).
\end{itemize}

Everything else is as without a project. All checks of \opt{--save} and
\opt{--resume} apply to the combined settings. The project file counts as a file that
\opt{--save} must not replace. Paths on the command line start from the current
folder, even when \opt{--project} names a file elsewhere. The \code{wrote} line shows
an absolute name when the output folder comes from the project file. If you give the
folder with \opt{-o}, the line shows it as you typed it.

\section{Safety}

When hdl-vet writes an output tree, it replaces the output folder. So the output folder,
wherever it comes from, must not contain a script of the run or the project file. The
\code{output} of a project must not be the folder of the project. It must not contain
this folder. Chapter~\ref{ch:cli} gives the rules for \opt{--force} on an output folder
from the project file.

\cmd{vet init} also writes these things:

\begin{itemize}
\item a \file{.gitignore} with \code{/out/} and \code{*.session.zip}
\item a \file{README.md} with the commands
\item a \file{cells} folder with a note about where to keep files for
  \apiname{Vet.add_file}
\end{itemize}

\lstinputlisting[style=console]{examples/console/project.txt}

The transcript shows these steps:

\begin{itemize}
\item \cmd{vet init} runs twice. The second time, the folder is not empty, and hdl-vet
  refuses.
\item A run in the project with an extra script. The script \file{skipped.py} runs
  after \file{01_example.py}, and the header line lists both.
\item \cmd{vet apply} writes \file{out}. hdl-vet refuses the same command again, because
  \file{out} exists. With \opt{--force}, hdl-vet replaces the folder. This works because
  the old \file{out} is an earlier output of hdl-vet with nothing else in it.
\item \opt{-f} on the command line. hdl-vet prints the note about the scripts of the
  project.
\item \opt{--project} from another folder, then the options that conflict, then a run
  without a project and without a file list.
\end{itemize}
